\documentclass[11pt,a4paper]{article}

% =========================================================================
%  RAPPORT INFRASTRUCTURE MULTI-SITES — RT Bank
%  Theme : Banque / infra entreprise — palette deep blue
% =========================================================================

\usepackage[utf8]{inputenc}
\usepackage[T1]{fontenc}
\usepackage[french]{babel}
\usepackage{lmodern}
\usepackage[a4paper,margin=2.2cm,top=2.6cm,bottom=2.4cm]{geometry}
\usepackage{xcolor}
\usepackage{fontawesome5}
\usepackage{tikz}
\usetikzlibrary{shapes.geometric,positioning,calc,arrows.meta}
\usepackage{titlesec}
\usepackage[most]{tcolorbox}
\usepackage{listings}
\usepackage{fancyhdr}
\usepackage{booktabs}
\usepackage{tabularx}
\usepackage{array}
\usepackage{colortbl}
\usepackage{enumitem}
\usepackage{microtype}
\usepackage[bookmarks=true,colorlinks=true,linkcolor=accent,urlcolor=accent,citecolor=accent]{hyperref}

% --- Palette banque / enterprise ----------------------------------------
\definecolor{primary}{HTML}{0F1F44}
\definecolor{accent}{HTML}{1E40AF}
\definecolor{accent2}{HTML}{D4A537}
\definecolor{accentdark}{HTML}{1E3A8A}
\definecolor{light}{HTML}{EFF6FF}
\definecolor{midgray}{HTML}{6B7280}
\definecolor{codebg}{HTML}{0F1A33}
\definecolor{codefg}{HTML}{E8EEF7}
\definecolor{codekw}{HTML}{93C5FD}
\definecolor{codestr}{HTML}{FBD38D}
\definecolor{codecmt}{HTML}{8FA3C7}

\titleformat{\section}{\sffamily\Large\bfseries\color{primary}}{\thesection}{0.8em}{}[\vspace{-0.4ex}{\color{accent}\rule{\linewidth}{1.2pt}}]
\titleformat{\subsection}{\sffamily\large\bfseries\color{accent}}{\thesubsection}{0.6em}{}
\titleformat{\subsubsection}{\sffamily\normalsize\bfseries\color{primary}}{\thesubsubsection}{0.5em}{}
\titlespacing*{\section}{0pt}{1.6em}{0.8em}

\pagestyle{fancy}
\fancyhf{}
\renewcommand{\headrulewidth}{0pt}
\fancyhead[L]{\footnotesize\sffamily\color{midgray}\faServer~RT Bank — Infra multi-sites sécurisée}
\fancyhead[R]{\footnotesize\sffamily\color{midgray}HONORINE Kylian}
\fancyfoot[C]{\footnotesize\sffamily\color{midgray}\thepage}

\lstdefinestyle{cfg}{
  backgroundcolor=\color{codebg},
  basicstyle=\ttfamily\footnotesize\color{codefg},
  keywordstyle=\color{codekw}\bfseries,
  stringstyle=\color{codestr},
  commentstyle=\color{codecmt}\itshape,
  showstringspaces=false, breaklines=true,
  frame=none, framesep=8pt, framerule=0pt,
  xleftmargin=12pt, xrightmargin=12pt,
  columns=fullflexible, inputencoding=utf8, extendedchars=true,
  morekeywords={zone,type,master,slave,allow-transfer,also-notify,file,directory,recursion,listen-on,allow-query,forwarders,frontend,backend,bind,server,acl,http-request,redirect,scheme,mode,balance,option,ssl,crt,verify,timeout,maxconn,daemon,global,defaults,log,VirtualHost,ServerName,DocumentRoot,SSLEngine,SSLCertificateFile,SSLCertificateKeyFile,Redirect,Order,allow,deny,Require,from,all,Listen,RewriteEngine,RewriteRule},
  literate={é}{{\'e}}1 {è}{{\`e}}1 {ê}{{\^e}}1 {à}{{\`a}}1 {ô}{{\^o}}1 {î}{{\^\i}}1 {ç}{{\c{c}}}1 {ù}{{\`u}}1 {â}{{\^a}}1 {É}{{\'E}}1 {È}{{\`E}}1 {Ê}{{\^E}}1 {À}{{\`A}}1 {Ô}{{\^O}}1 {Ç}{{\c{C}}}1 {Ù}{{\`U}}1 {Â}{{\^A}}1
}
\lstset{style=cfg}

\newtcblisting{cfgbox}[1][]{
  listing only,
  enhanced jigsaw, breakable,
  colback=codebg, colframe=accent,
  arc=2pt, boxrule=0pt, leftrule=3pt,
  left=8pt, right=8pt, top=4pt, bottom=4pt,
  listing options={style=cfg}, #1
}

\newtcolorbox{infobox}[1][]{
  enhanced, colback=light, colframe=accent,
  arc=2pt, boxrule=0pt, leftrule=3pt,
  left=10pt, right=10pt, top=6pt, bottom=6pt,
  fontupper=\small, #1
}
\newtcolorbox{secbox}[1][]{
  enhanced, colback=accent!8, colframe=accent2,
  arc=2pt, boxrule=0pt, leftrule=3pt,
  left=10pt, right=10pt, top=6pt, bottom=6pt,
  fontupper=\small, #1
}
\newtcolorbox{alertbox}[1][]{
  enhanced, colback=red!5, colframe=red!70!black,
  arc=2pt, boxrule=0pt, leftrule=3pt,
  left=10pt, right=10pt, top=6pt, bottom=6pt,
  fontupper=\small, #1
}

\setlist[itemize]{leftmargin=*,itemsep=2pt,topsep=4pt}
\setlist[enumerate]{leftmargin=*,itemsep=2pt,topsep=4pt}

\begin{document}
\pagenumbering{gobble}

\begin{tikzpicture}[remember picture, overlay]
  \fill[primary] (current page.north west) rectangle ([yshift=-9cm]current page.north east);
  \fill[accent2] ([yshift=-9cm]current page.north west) rectangle ([yshift=-9.25cm]current page.north east);
  % Logo : building / bank avec cadenas
  \node[anchor=center] at ([yshift=-4.6cm]current page.north) {%
    \begin{tikzpicture}[scale=1.0]
      % toit en triangle
      \fill[white] (-2.4,0.6) -- (0,1.7) -- (2.4,0.6) -- cycle;
      % base
      \fill[white] (-2.2,-1.4) rectangle (2.2,0.6);
      % colonnes
      \foreach \x in {-1.6,-0.8,0,0.8,1.6} {
        \fill[accent2] (\x-0.18,-1.2) rectangle (\x+0.18,0.55);
      }
      % escalier
      \fill[white] (-2.5,-1.6) rectangle (2.5,-1.4);
      \fill[white] (-2.7,-1.85) rectangle (2.7,-1.6);
      % colonne centrale : icone serveur stack
      \fill[accent] (-0.2,-0.6) rectangle (0.2,0.4);
      % toit details
      \fill[accent2] (-0.15,1.0) rectangle (0.15,1.55);
      % petits cadenas qui flottent autour
      \node[scale=0.9,accent2] at (-2.8,1.4) {\faLock};
      \node[scale=0.7,accent2!70] at (2.7,1.6) {\faLock};
      \node[scale=0.6,white!70!primary] at (-3.0,-0.6) {\faShieldVirus};
      \node[scale=0.6,white!70!primary] at (3.0,-0.4) {\faShieldVirus};
    \end{tikzpicture}};
  \node[anchor=north, white, font=\sffamily\Huge\bfseries] at ([yshift=-6.7cm]current page.north) {RT Bank};
  \node[anchor=north, accent2, font=\sffamily\Large] at ([yshift=-7.65cm]current page.north) {Infrastructure multi-sites sécurisée};
  \node[anchor=north, white!75!primary, font=\sffamily\small] at ([yshift=-8.4cm]current page.north) {Stormshield \textbullet{} Apache HA \textbullet{} BIND9 \textbullet{} AD \textbullet{} Zabbix};

  \node[anchor=south west, primary, font=\sffamily\small] at ([xshift=2.2cm,yshift=3cm]current page.south west) {\textbf{\textcolor{accent}{Auteur}}};
  \node[anchor=south west, primary, font=\sffamily\small] at ([xshift=2.2cm,yshift=2.5cm]current page.south west) {HONORINE Kylian};
  \node[anchor=south east, primary, font=\sffamily\small] at ([xshift=-2.2cm,yshift=3cm]current page.south east) {\textbf{\textcolor{accent}{Cadre}}};
  \node[anchor=south east, primary, font=\sffamily\small] at ([xshift=-2.2cm,yshift=2.5cm]current page.south east) {SAE 4.Cyber.04};
  \fill[accent2] ([yshift=1.3cm]current page.south west) rectangle ([yshift=1.4cm]current page.south east);
  \node[anchor=south, midgray, font=\sffamily\footnotesize] at ([yshift=0.6cm]current page.south) {IUT Réseaux \& Télécommunications — Filière Cybersécurité};
\end{tikzpicture}
\newpage

\pagenumbering{arabic}
\setcounter{page}{1}

{\sffamily
\hypersetup{linkcolor=primary}
\renewcommand{\contentsname}{\textcolor{accent}{\faList~} Sommaire}
\tableofcontents
}
\vspace{0.4cm}

\section{Contexte et objectifs}

\begin{infobox}[title={\faBullseye~Mission}]
RT Bank \textemdash{} agence bancaire fictive \textemdash{} déploie une infrastructure multi-sites entre Saint-Denis et Saint-Pierre. Objectifs : sécurité par défaut, haute disponibilité, exploitation supervisée, et conformité aux recommandations \textbf{ANSSI}.
\end{infobox}

\textbf{Périmètre couvert :}
\begin{itemize}
  \item Pare-feu \textbf{Stormshield EVA} segmentant DMZ / LAN-Clients / LAN-Serveurs / WAN
  \item Cluster Apache + HAProxy avec terminaison TLS en DMZ
  \item DNS \textbf{BIND9} maître/esclave pour le domaine \texttt{rt-bank.re}
  \item \textbf{Active Directory} \texttt{rtbank.re} avec DC primaire à Saint-Denis et RODC à Saint-Pierre
  \item Supervision \textbf{Zabbix} et durcissement système
  \item Pentest et plan de contre-mesures
\end{itemize}

% =========================================================================
\section{Architecture réseau}

\subsection{Plan d'adressage}

\renewcommand{\arraystretch}{1.3}
\begin{tabularx}{\linewidth}{@{}>{\bfseries}l l l X@{}}
\toprule
\rowcolor{light} Zone & VLAN & Réseau & Rôle \\ \midrule
WAN           & --  & 192.168.41.0/23 & Sortie Internet \\
DMZ           & 100 & 10.10.10.0/24   & Services exposés (web, mail relay) \\
LAN-Serveurs  & 200 & 10.10.20.0/24   & DC, DNS, fichiers, Zabbix \\
LAN-Clients   & 300 & 10.10.30.0/24   & Postes utilisateurs \\
MGMT          & 999 & 10.10.99.0/24   & Administration équipements \\
\bottomrule
\end{tabularx}

\subsection{Stormshield EVA — interfaces}

\begin{tabularx}{\linewidth}{@{}c >{\bfseries}l l X@{}}
\toprule
\rowcolor{light} Iface & Zone & IP & Description \\ \midrule
1 & WAN          & 192.168.41.65/23 & Lien opérateur \\
2 & DMZ          & 10.10.10.1/24    & Frontaux web \\
3 & LAN-Clients  & 10.10.30.1/24    & Postes \\
4 & LAN-Serveurs & 10.10.20.1/24    & Backoffice \\
\bottomrule
\end{tabularx}

\subsection{Matrice de flux}

\renewcommand{\arraystretch}{1.4}
\begin{tabularx}{\linewidth}{@{}l l l c X@{}}
\toprule
\rowcolor{light} Source & Destination & Service & Action & Justification \\ \midrule
LAN-Clients & DMZ (HAProxy)  & TCP/443         & \textcolor{accent}{\faCheck}~OK  & Accès web interne \\
WAN         & DMZ (HAProxy)  & TCP/443         & \textcolor{accent}{\faCheck}~OK  & Web public \\
LAN-Clients & LAN-Serveurs (DC) & TCP/389+88 + UDP/53 & \textcolor{accent}{\faCheck}~OK & Auth + DNS \\
DMZ         & LAN-Serveurs   & TCP/3306        & \textcolor{accent}{\faCheck}~OK  & SQL applicative \\
DMZ         & LAN-Clients    & Any             & \textcolor{red}{\faTimes}~Refus  & Isolation stricte \\
WAN         & LAN-*          & Any             & \textcolor{red}{\faTimes}~Refus  & Aucun accès direct \\
\bottomrule
\end{tabularx}

% =========================================================================
\section{Frontaux web — HAProxy + Apache + TLS}

\subsection{Schéma fonctionnel}

\begin{secbox}[title={\faProjectDiagram~Chaîne de traitement}]
\textbf{Internet} $\rightarrow$ \textbf{Stormshield (NAT 443)} $\rightarrow$ \textbf{HAProxy (terminaison TLS, équilibrage)} $\rightarrow$ \textbf{2 Apache backends} (\texttt{web1}, \texttt{web2})
\end{secbox}

HAProxy joue trois rôles : terminer le TLS (1 seul certificat à gérer), répartir la charge (round-robin avec health checks), et masquer le nombre réel de frontaux.

\subsection{Configuration HAProxy}

\begin{cfgbox}
global
    daemon
    maxconn 4096
    log /dev/log local0

defaults
    mode http
    timeout connect 5s
    timeout client  30s
    timeout server  30s
    option httplog
    option forwardfor

frontend https_in
    bind *:443 ssl crt /etc/haproxy/certs/rtbank.pem
    http-request redirect scheme https unless { ssl_fc }
    default_backend web_pool

backend web_pool
    balance roundrobin
    option httpchk GET /health
    server web1 10.10.10.11:80 check
    server web2 10.10.10.12:80 check
\end{cfgbox}

\subsection{Apache backend (web1 \& web2)}

Apache écoute sur le port 80 en interne uniquement. La redirection \texttt{X-Forwarded-Proto} est lue depuis HAProxy pour générer correctement les liens absolus.

\begin{cfgbox}
<VirtualHost *:80>
    ServerName www.rt-bank.re
    DocumentRoot /var/www/rtbank

    # Endpoint sante pour HAProxy
    Alias /health /var/www/health
    <Directory /var/www/health>
        Require all granted
    </Directory>

    <Directory /var/www/rtbank>
        Options -Indexes +FollowSymLinks
        AllowOverride None
        Require all granted
    </Directory>

    ErrorLog  ${APACHE_LOG_DIR}/rtbank-error.log
    CustomLog ${APACHE_LOG_DIR}/rtbank-access.log combined
</VirtualHost>
\end{cfgbox}

% =========================================================================
\section{DNS — BIND9 maître / esclave}

\subsection{Rôles}

\renewcommand{\arraystretch}{1.3}
\begin{tabularx}{\linewidth}{@{}>{\bfseries}l l X@{}}
\toprule
\rowcolor{light} Serveur & Adresse & Rôle \\ \midrule
\texttt{ns1.rt-bank.re} & 10.10.20.5 & Master \texttt{rt-bank.re} (Saint-Denis) \\
\texttt{ns2.rt-bank.re} & 10.10.20.6 & Slave  \texttt{rt-bank.re} (Saint-Pierre) \\
\bottomrule
\end{tabularx}

\subsection{Configuration du serveur maître}

\begin{cfgbox}
options {
    directory "/var/cache/bind";
    listen-on { 127.0.0.1; 10.10.20.5; };
    allow-query { 10.10.0.0/16; };
    forwarders { 1.1.1.1; 9.9.9.9; };
    recursion yes;
    dnssec-validation auto;
};

zone "rt-bank.re" {
    type master;
    file "/etc/bind/db.rt-bank.re";
    allow-transfer { 10.10.20.6; };
    also-notify    { 10.10.20.6; };
};
\end{cfgbox}

\subsection{Zone \texttt{rt-bank.re}}

\begin{cfgbox}
$TTL 3600
@   IN SOA ns1.rt-bank.re. admin.rt-bank.re. (
        2026053001  ; serial
        3600        ; refresh
        900         ; retry
        604800      ; expire
        3600 )      ; minimum

    IN NS    ns1.rt-bank.re.
    IN NS    ns2.rt-bank.re.

ns1 IN A     10.10.20.5
ns2 IN A     10.10.20.6
www IN A     192.168.41.66      ; IP publique HAProxy
mail IN A    192.168.41.67
@   IN MX 10 mail.rt-bank.re.
\end{cfgbox}

\subsection{Configuration du serveur esclave}

\begin{cfgbox}
zone "rt-bank.re" {
    type slave;
    masters { 10.10.20.5; };
    file "/var/cache/bind/db.rt-bank.re.slave";
};
\end{cfgbox}

% =========================================================================
\section{Annuaire — Active Directory \texttt{rtbank.re}}

\subsection{Topologie}

\begin{secbox}[title={\faNetworkWired~Réplication AD}]
\begin{itemize}
  \item \textbf{DC1 — Saint-Denis} (10.10.20.10) : DC principal, FSMO roles, schema master
  \item \textbf{RODC — Saint-Pierre} (10.10.20.11) : Read-Only Domain Controller pour le site distant
\end{itemize}
\end{secbox}

Un \textbf{RODC} (Read-Only DC) authentifie localement les utilisateurs autorisés sans pouvoir être source de propagation en cas de compromission \textemdash{} idéal pour un site distant moins protégé physiquement.

\subsection{Stratégies appliquées}

\renewcommand{\arraystretch}{1.3}
\begin{tabularx}{\linewidth}{@{}>{\bfseries}l X@{}}
\toprule
\rowcolor{light} GPO & Effet \\ \midrule
Password policy   & 14 caractères, complexité, expiration 90j, historique 24 \\
Lockout policy    & 5 tentatives, 30 min de verrouillage \\
Audit policy      & Logon/Logoff, Account Mgmt, Policy Change \textemdash{} tous succès/échecs \\
UAC               & Niveau maximum sur les postes utilisateurs \\
LAPS              & Mot de passe admin local unique par poste, rotation 30j \\
\bottomrule
\end{tabularx}

\subsection{Structuration des OUs}

\begin{cfgbox}[listing options={style=cfg,language=bash}]
DC=rtbank,DC=re
|-- OU=Sites
|   |-- OU=StDenis
|   |   |-- OU=Users      (collaborateurs)
|   |   |-- OU=Computers  (postes)
|   |   `-- OU=Servers    (serveurs)
|   `-- OU=StPierre
|       |-- OU=Users
|       `-- OU=Computers
|-- OU=Groups
|   |-- OU=Roles       (RH, Compta, Direction, IT)
|   `-- OU=Resources   (acces partages, applis)
`-- OU=ServiceAccounts (comptes applicatifs)
\end{cfgbox}

% =========================================================================
\section{Supervision — Zabbix}

\subsection{Architecture}

\begin{tabularx}{\linewidth}{@{}>{\bfseries}l X@{}}
\toprule
\rowcolor{light} Composant & Description \\ \midrule
Zabbix server    & 10.10.20.20 — collecte centralisée \\
MySQL backend    & sur le même serveur, base \texttt{zabbix} dédiée \\
Frontend         & Apache + PHP, accès via HAProxy interne \\
Agents Zabbix    & sur DC, DNS, frontaux web, pare-feu (via SNMPv3) \\
\bottomrule
\end{tabularx}

\subsection{Templates et alertes clés}

\begin{tabularx}{\linewidth}{@{}>{\bfseries}p{4cm} X@{}}
\toprule
\rowcolor{light} Métrique & Seuil d'alerte \\ \midrule
CPU utilisation  & \textgreater{}~85\% pendant 5 min \\
RAM disponible   & \textless{}~10\% pendant 5 min \\
Disque \texttt{/}, \texttt{/var}, \texttt{/srv} & \textgreater{}~90\% \\
Process critiques (apache2, named, haproxy) & Indisponible \textgreater{}~60\,s \\
Réplication AD   & SYSVOL/NETLOGON non répliqué \textgreater{}~15 min \\
Certificat TLS   & Expiration \textless{}~30 jours \\
\bottomrule
\end{tabularx}

\subsection{Notifications}

Trois canaux : mail SMTP via le relais DMZ pour les alertes \og Information \fg{} et \og Avertissement \fg{}, webhook Mattermost pour les alertes \og Haute \fg{}, et SMS via passerelle dédiée pour les alertes \og Critique \fg{}.

% =========================================================================
\section{Durcissement — recommandations ANSSI}

\subsection{Systèmes Linux (Debian 12)}

\begin{cfgbox}[listing options={style=cfg,language=bash}]
# Service indispensable au strict minimum
systemctl disable --now bluetooth cups avahi-daemon

# Pare-feu hote (UFW) ferme par defaut
ufw default deny incoming
ufw default allow outgoing
ufw allow from 10.10.99.0/24 to any port 22 proto tcp
ufw enable

# Mises a jour automatiques securite
apt install -y unattended-upgrades
dpkg-reconfigure --priority=low unattended-upgrades

# SSH durci
sed -i 's/^#\?PermitRootLogin.*/PermitRootLogin no/' /etc/ssh/sshd_config
sed -i 's/^#\?PasswordAuthentication.*/PasswordAuthentication no/' /etc/ssh/sshd_config
sed -i 's/^#\?X11Forwarding.*/X11Forwarding no/' /etc/ssh/sshd_config
systemctl restart sshd

# Auditd active
apt install -y auditd
systemctl enable --now auditd
\end{cfgbox}

\subsection{Windows Server (DC1)}

\begin{itemize}
  \item Désactivation de SMBv1 (\texttt{Set-SmbServerConfiguration -EnableSMB1Protocol \$false})
  \item Activation du \textbf{Credential Guard} et \textbf{LSA Protection}
  \item Blocage NTLMv1 / NTLMv2 fallback via GPO
  \item Activation Windows Defender ASR (Attack Surface Reduction)
  \item Journaux exportés vers le serveur Zabbix + retention 1 an
\end{itemize}

% =========================================================================
\section{Tests d'intrusion et contre-mesures}

\subsection{Méthodologie}

Trois phases : reconnaissance, exploitation, post-exploitation. Outils : nmap, gobuster, hydra, BloodHound, sqlmap. Périmètre : DMZ et services exposés.

\subsection{Findings principaux et remédiations}

\renewcommand{\arraystretch}{1.4}
\begin{tabularx}{\linewidth}{@{}>{\bfseries}p{4cm} X@{}}
\toprule
\rowcolor{light} Vulnérabilité & Contre-mesure \\ \midrule
Énumération via NSE (\texttt{nmap -sV})       & Filtrage applicatif sur le pare-feu, masquage de version Apache (\texttt{ServerTokens Prod}) \\
Bruteforce HTTP form               & Rate-limiting HAProxy + fail2ban sur les logs Apache \\
Injection SQL (sqlmap)             & Paramétrage des requêtes en backend, WAF ModSecurity activé en mode bloquant \\
Reconnaissance AD (BloodHound)     & Tiering AD (Tier 0/1/2), suppression des chemins d'attaque vers DA \\
Kerberoasting / AS-REP roasting    & Mots de passe service \textgreater{}~25 caractères, gMSA quand possible, \texttt{Do not require pre-authentication} désactivé \\
Mouvement latéral SMB              & Désactivation SMBv1, signature SMB obligatoire, segmentation par VLAN \\
\bottomrule
\end{tabularx}

\begin{alertbox}[title={\faExclamationTriangle~Constat majeur}]
Le risque résiduel le plus élevé porte sur \textbf{le facteur humain} \textemdash{} phishing ciblé, partage de mots de passe, branchement de matériel non maîtrisé. Une démarche de sensibilisation continue et des exercices red team réguliers sont indispensables.
\end{alertbox}

% =========================================================================
\section{Synthèse}

\renewcommand{\arraystretch}{1.35}
\begin{tabularx}{\linewidth}{@{}>{\bfseries}p{3.5cm} X@{}}
\toprule
\rowcolor{light} Volet & Livrables \\ \midrule
Réseau      & Plan d'adressage, segmentation VLAN, matrice de flux \\
Sécurité    & Stormshield (filtrage + NAT), durcissement ANSSI \\
Services    & HAProxy/Apache HA + TLS, BIND9 master/slave, AD multi-sites \\
Supervision & Templates Zabbix, escalade multi-canaux \\
Pentest     & Rapport \textit{findings} + plan de remédiation \\
\bottomrule
\end{tabularx}

\section{Conclusion}

Cette SAE met en \oe{}uvre une infrastructure d'entreprise complète, depuis le filtrage en bordure jusqu'à la supervision applicative. La logique est invariante : \textbf{minimum de privilèges}, \textbf{défense en profondeur}, \textbf{traçabilité}. Aucun composant n'est seul : chaque service critique (DNS, AD, web) est redondé, chaque flux passe par un point de contrôle, et chaque incident est instrumenté pour être détecté en moins d'une minute.

\medskip
L'expérience la plus instructive a été le pentest : voir le réseau du point de vue de l'attaquant, mesurer l'écart entre la documentation et la réalité, et constater que le maillon faible est presque toujours organisationnel avant d'être technique.

\vspace{0.5cm}
\begin{center}
{\sffamily\itshape\color{midgray}\large
\og La sécurité, c'est un processus \textemdash{} pas un produit. \fg{} \\ \small{— Bruce Schneier}}
\end{center}

\end{document}
